\documentclass[dvipdfmx]{article}
\setcounter{secnumdepth}{5}
\setcounter{tocdepth}{1}

\usepackage{amsmath}
\usepackage{amssymb}
\usepackage{amsthm}
\usepackage{enumerate}
\usepackage{tikz}
\usetikzlibrary{cd}
\usepackage{mathtools}


\usepackage[dvipdfmx]{hyperref}
\usepackage{pxjahyper}
\hypersetup{
 setpagesize=false,
 bookmarksnumbered=true,%
 bookmarksopen=true,%
 colorlinks=true,%
 linkcolor=blue,
 citecolor=red,
}

\theoremstyle{definition}

\newtheorem{thm}{Theorem}[section]
\newtheorem{prop}[thm]{Proposition}
\newtheorem{lem}[thm]{Lemma}
\newtheorem{defn}[thm]{Definition}
\newtheorem{rem}[thm]{Remark}
\usepackage[utf8]{inputenc}


\renewcommand{\emph}{\textbf}

\usepackage{enumitem}
\setlist[enumerate,1]{label=(\roman*)}


\title{20260513}
\author{Kazuki Masugi}
\date{}
\begin{document}
\maketitle

\section{ab theory}

Albanese 版の new ordinary のような理論があってもいいのではないか ? というより、polycurve での ab theory について、よくわからないとは思っているが、そもそも ($\ell$ での挙動を多少でも有益なかたちで制御できないか ?)

\section{memo}

いくつかノートをつくりたい話題がたまっている。

\begin{itemize}
	\item local object についての話
	\item log 的な仕方で、群論的な Chern 類をみる話
	\item log 基本群まわりの基本的な話、特に exact sequence まわり
\end{itemize}









\end{document}

















